\section*{Chapitre \ref{ch:g11:absorbance} --- Couleur et absorbance}

\begin{solution}{exo:g11:absorbance:1}
Si elle absorbe l'orange, elle paraît bleue, la couleur opposée à
l'orange sur le cercle. Si elle absorbe le violet, elle paraît jaune.
\end{solution}

\begin{solution}{exo:g11:absorbance:2}
Une solution verte absorbe principalement la lumière rouge, la couleur
opposée au vert sur le cercle~: d'environ \qty{620}{nm} à
\qty{750}{nm}.
\end{solution}

\begin{solution}{exo:g11:absorbance:3}
Colorant bleu~: $\lambda_{\max} = \qty{629}{nm}$, $A = 0.82$. Colorant
jaune~: $\lambda_{\max} = \qty{427}{nm}$, $A = 0.795$, soit environ 0.80.
\end{solution}

\begin{solution}{exo:g11:absorbance:4}
$A$ est proportionnelle à la concentration~: trois fois plus donne
$A = 0.90$~; deux fois moins donne $A = 0.15$.
\end{solution}

\begin{solution}{exo:g11:absorbance:5}
Le blanc est une cuve remplie du \omterm{def:g6:solutions-and-solubility:solute}{solvant} pur. Faire $A = 0$ avec lui
élimine ce qu'absorbent la cuve, le \omterm{def:g6:solutions-and-solubility:solute}{solvant} et l'instrument, si bien que
les mesures suivantes sont dues au seul \omterm{def:g6:solutions-and-solubility:solute}{soluté}.
\end{solution}

\begin{solution}{exo:g11:absorbance:6}
$C = 0.41 / 0.164 = \qty{2.5}{mg/L}$.
\end{solution}

\begin{solution}{exo:g11:absorbance:7}
À \qty{629}{nm}, l'\omterm{def:g11:absorbance:absorbance}{absorbance} est la plus grande~: la mesure est la plus
sensible (de petites concentrations donnent encore des \omterm{def:g11:absorbance:absorbance}{absorbances}
lisibles)~; et au sommet de la bande la courbe est plate, si bien qu'une
petite erreur sur la longueur d'onde change à peine la mesure. À
\qty{550}{nm}, l'\omterm{def:g11:absorbance:absorbance}{absorbance} est faible et varie vite avec la longueur
d'onde.
\end{solution}

\begin{solution}{exo:g11:absorbance:8}
Diluée~: $C = 0.48 / 0.164 = \qty{2.9}{mg/L}$~; la boisson~:
$5 \times 2.9 = \qty{15}{mg/L}$ (plus précisément $14.6$).
\end{solution}

\begin{solution}{exo:g11:absorbance:9}
$a = A / (\ell\, C_m) = 0.795 / (1.00 \times 0.0150) = \qty{53.0}{L/(g.cm)}$;
$\varepsilon = a \times M = 53.0 \times 534.4 = \qty{2.83e4}{L/(mol.cm)}$.
\end{solution}

\begin{solution}{exo:g11:absorbance:10}
Au-dessus d'environ \qty{520}{nm}. À \qty{629}{nm}, le colorant jaune
n'absorbe pas du tout~: toute l'\omterm{def:g11:absorbance:absorbance}{absorbance} y est due au colorant bleu.
\end{solution}

\begin{solution}{exo:g11:absorbance:11}
Elle double~: $A = \varepsilon \ell c$ est proportionnelle à $\ell$~; la
lumière traverse deux fois plus de solution, donc deux fois plus de
\omterm{def:g7:atoms-and-molecules:molecule}{molécules} absorbantes.
\end{solution}

\begin{solution}{exo:g11:absorbance:12}
$A/C$~: $0.164$, $0.162$, $0.115$, $0.0675$. Le rapport n'est constant
que pour les deux premiers~: au-delà d'environ \qty{10}{mg/L}
(\omterm{def:g11:absorbance:absorbance}{absorbances} supérieures à 1.6 environ), la loi n'est plus vérifiée.
Seuls les étalons à 5 et \qty{10}{mg/L} (et les plus dilués) peuvent
servir~; une inconnue qui donne 2.5 doit être diluée, par exemple cinq
fois, puis mesurée de nouveau.
\end{solution}

\begin{solution}{exo:g11:absorbance:13}
Colorant bleu, d'après \qty{629}{nm} seule~: $C = 0.574 / 0.164 = \qty{3.5}{mg/L}$.
Colorant jaune, d'après \qty{427}{nm}~: $C = 0.53 / 0.0530 = \qty{10}{mg/L}$.
\end{solution}

\begin{solution}{exo:g11:absorbance:14}
Elle trouve $0.368 / 0.164 = \qty{2.24}{mg/L}$. La vraie \omterm{def:g11:absorbance:absorbance}{absorbance} est
$0.368 - 0.040 = 0.328$, si bien que la vraie concentration vaut
$0.328 / 0.164 = \qty{2.00}{mg/L}$. Elle se trompe de
\qty{12}{\percent} par excès.
\end{solution}

\begin{solution}{exo:g11:absorbance:15}
$10 \times 60 = \qty{600}{mg}$ par jour, dans $600 / 20 = \qty{30}{L}$ de
boisson. Personne ne boit trente litres par jour~: la boisson seule ne
présente pas de risque réaliste, même si le colorant peut aussi venir
d'autres aliments.
\end{solution}

\begin{solution}{pb:g11:absorbance:1}
\textbf{1.} Le rouge-orangé, la couleur opposée au bleu sur le cercle.

\textbf{2.} $\lambda_{\max} = \qty{629}{nm}$.

\textbf{3.} À \qty{450}{nm}, la lumière est bleue~: le colorant la
laisse passer, et c'est exactement pour cela que la boisson paraît bleue.

\textbf{4.} À \qty{629}{nm}.

\textbf{5.} $\qty{50.0}{mg} / \qty{0.5000}{L} = \qty{100}{mg/L}$.

\textbf{6.} \omterm{def:g10:concentration-and-dilution:dilution}{Facteur de dilution} $100 / 3.0 = 33.3$, donc
$V_0 = 100.0 / 33.3 = \qty{3.00}{mL}$~: une pipette graduée (ou une
burette) et une fiole jaugée de \qty{100.0}{mL}.

\textbf{7.} $A/C$ = 0.168, 0.163, 0.165, 0.163, 0.165~: toutes proches
de 0.164, donc $A \approx 0.164\, C$, une droite passant par l'origine.

\textbf{8.} Une solution sans colorant ($C = 0$) est le blanc, réglé à
$A = 0$.

\textbf{9.} $a = \qty{0.164}{L/mg}$ par centimètre, soit
\qty{164}{L/(g.cm)}~; $\varepsilon = 164 \times 792.9 =
\qty{1.30e5}{L/(mol.cm)}$.

\textbf{10.} $F = 50.0 / 25.0 = 2$.

\textbf{11.} $C = 0.590 / 0.164 = \qty{3.60}{mg/L}$.

\textbf{12.} $2 \times 3.60 = \qty{7.20}{mg/L}$.

\textbf{13.} Environ $2 \times 0.590 = 1.18$, au-dessus de l'étalon le
plus concentré (0.823)~: la mesure sortirait du domaine étalonné, où la
loi peut ne plus être vérifiée. La \omterm{def:g10:concentration-and-dilution:dilution}{dilution} la ramène parmi les
étalons.

\textbf{14.} $\qty{7.20}{mg/L} \times \qty{0.500}{L} = \qty{3.60}{mg}$.

\textbf{15.} $n = \num{3.60e-3} / 792.9 = \qty{4.54e-6}{mol}$.

\textbf{16.} $6 \times 30 = \qty{180}{mg}$ par jour.

\textbf{17.} $180 / 3.60 = 50$ bouteilles.

\textbf{18.} $50 \times 0.500 = \qty{25}{L}$ en une journée~: impossible
à boire. Le colorant de cette boisson ne peut pas, à lui seul, amener un
enfant près de la limite, même si les colorants de plusieurs aliments
s'additionnent.

\textbf{19.} $6 \times 70 = \qty{420}{mg}$, soit $420 / 3.60 \approx
117$ bouteilles.

\textbf{20.} Environ 50 bouteilles de \qty{500}{mL}, soit \qty{25}{L} de
boisson, en une seule journée.
\end{solution}
